\section{Discussion and conclusion}

\subsection{What the controls change}

The experiments distinguish concentration in activations, selection within a candidate grid, and an economic factor count. In the original diagnostic family, rankings change across assets, weighting, archived pipelines, seeds, and periods. The matched controls narrow the explanation: absolute rank also falls in random networks as layers narrow, training adds concentration beyond that benchmark, and removing six features on a common panel does not reproduce the archived pipeline reversal. The original contrast cannot be assigned to feature deletion alone.

The conditional-autoencoder extension changes the scope of the selection result. Its point-estimate diagnostics are favorable relative to the chosen reference, and selection is concentrated at $K=8$ in the fixed grid. The study therefore contains both sensitive and comparatively concentrated selections. It does not establish universal dimension instability among successful pricing models. Because eight is the largest candidate, the result also does not identify a global optimal width. A finite-grid preference and the number of priced structural shocks remain different objects even when that preference is stable.

Inference further limits the interpretation. All 15 exploratory pairwise intervals include zero, but the primary known-truth calibration exhibits undercoverage. We retain that result and do not use non-rejection as evidence of equivalence or structural nonidentification. The failure is specific to the tested procedure and moment designs; it does not estimate empirical error probabilities or invalidate every other procedure in the study. Nor is it an economic explanation for why a dimension changes. Its role is to delimit the inferential claims available from the extension.

\subsection{Economic interpretation and reporting}

The original market-direction diagnostics suggest a distinction between dominant common payoff variation and relative-pricing errors. Removing a validation-estimated market component attenuates the geometry contrast in later-period point estimates, while adding the market factor does not jointly improve raw moments and alphas. Wide attenuation intervals and pipeline sensitivity prevent precise causal attribution of compression losses to market risk. We do not estimate the separate contributions of state complexity, risk exposures, or risk prices.

For empirical reporting, a selected width should be accompanied by the candidate loss profile, its payoff span and weighting, and algorithmic dispersion. A stated tolerance can summarize small observed gaps, but cannot establish equivalence without a justified margin and valid inference. Losses computed using evaluation-estimated weights in the original diagnostic exercise and validation-fixed weights in the conditional-autoencoder extension should not be combined into a model league table. Each comparison is interpreted within its own fixed evaluation protocol.

The separate temporal replay illustrates a related boundary. Retrospective forecast changes can align with subsequent errors without delivering a successful lagged switching rule or significant primary incremental portfolio gains. Because those predictors differ from the pricing networks and the universe was retrospectively selected, this exercise does not show that latent pricing dimension varies through time or that a prospective investor could earn the reported returns.

\subsection{Limitations}

The original pricing family failed earlier benchmark-advancement criteria. The later conditional-autoencoder control improves benchmark relevance by the prespecified point-estimate diagnostic, but does not establish state-of-the-art performance, superiority to its matched linear-loading arm, or a representative result for all neural pricing architectures. All results concern U.S. equities, conventional numerical characteristics, and the specified asset spans.

Chronology restricts confirmation claims. Legacy development included January 2020 before the 2020--2025 exercise was locked. The feature, architecture, conditional-autoencoder, and inference controls were designed after inspection of earlier findings. Their frozen run protocols limit further search but do not make the historical data new. The conditional-autoencoder extension does not access 2020--2025. None of these controls is described as an independent prospective replication.

The characteristic bridge retains benchmark values where available and reconstructed values otherwise. It is not an independent point-in-time vintage archive. Accounting revisions, coverage, and timing may affect rankings. The matched feature control preserves 184 input positions, rather than recreating the original 172-input architecture; its conclusions are conditional on that parameterization. Licensed security-level data cannot be included in the public package, so aggregate exhibit replay and licensed-data reconstruction are separate reproducibility tasks.

Geometry is measured at a finite cross-sectional sample and neighborhood scale. The tangent-angle diagnostic does not identify a unique smooth pricing manifold. Random-network comparisons use five fresh initialization draws rather than recovered original initial tensors; the functional probe contains 30 models and five seed folds. Its failed threshold does not prove that geometry contains no economic information. The known-loading simulation and penalty calculation concern a finite nested Gaussian family, not separately fitted neural networks. The new inference calibration conditions on fixed models and weights and does not reproduce uncertainty from the full training pipeline.

\subsection{Conclusion}

The evidence supports a conditional account of low-dimensional pricing representations. Training-associated concentration survives an architecture control, but the specified functional probe does not establish its incremental pricing role. The archived pipeline reversal does not survive interpretation as a matched six-feature deletion effect. A conditional autoencoder with favorable point-estimate diagnostics concentrates candidate selection at the upper grid boundary, qualifying the original instability narrative. Known-truth exercises distinguish selection error from economic dimension and expose a limitation of the new simultaneous intervals. These results identify which interpretations the tested comparisons support; they neither recover an invariant structural risk count nor prove that such a count cannot be identified by another design.
